% !TeX spellcheck = ru_RU
% !TEX root = vkr.tex

\section{Методология экспериментов} \label{sec:method}

В данном разделе описывается методология экспериментальной части исследования: пайплайн дообучения и оценки моделей, состав бенчмарков и обучающих датасетов, протокол дообучения, а также статистический дизайн эксперимента.

\subsection{Пайплайн дообучения и оценки} \label{subsec:pipeline}

Экспериментальный процесс организован в виде шестиэтапного пайплайна, обеспечивающего воспроизводимость и систематичность проведения экспериментов. Общая схема пайплайна представлена на Рис.~\ref{fig:pipeline}.

\begin{figure}[H]
\centering
\begin{tikzpicture}[
    node distance=1.2cm and 0.5cm,
    block/.style={
        rectangle, draw, rounded corners,
        text width=3.8cm, minimum height=1.1cm,
        align=center, font=\small
    },
    arrow/.style={->, thick, >=stealth},
]
    % Верхний ряд
    \node[block] (gen) {1. Генерация\\бенчмарка};
    \node[block, right=of gen] (prep) {2. Подготовка\\обучающих данных};
    \node[block, right=of prep] (train) {3. Дообучение\\(LoRA)};

    % Нижний ряд (справа налево)
    \node[block, below=of train] (infer) {4. Инференс\\на бенчмарке};
    \node[block, below=of prep] (eval) {5. Оценка\\и конвертация};
    \node[block, below=of gen] (stat) {6. Статистический\\анализ};

    % Стрелки
    \draw[arrow] (gen) -- (prep);
    \draw[arrow] (prep) -- (train);
    \draw[arrow] (train) -- (infer);
    \draw[arrow] (infer) -- (eval);
    \draw[arrow] (eval) -- (stat);
\end{tikzpicture}
\caption{Пайплайн дообучения и оценки моделей.}
\label{fig:pipeline}
\end{figure}

\noindent Каждый этап пайплайна выполняет следующие функции:
\begin{enumerate}
    \item \textbf{Генерация бенчмарка.} Формирование оценочных наборов данных двух типов: смешанного (реальные и синтетические запросы) и реального (только пользовательские запросы). Каждый пример содержит пользовательский запрос и эталонную метку инструмента (GEO или TAG).
    \item \textbf{Подготовка обучающих данных.} Создание обучающих датасетов различного объёма, типа и класcового баланса. Синтетические данные генерируются по методологии Self-Instruct~\cite{wang2023selfinstruct} с последующей гуманизацией для приближения к стилистике реальных запросов.
    \item \textbf{Дообучение (LoRA).} Адаптация базовой модели YandexGPT Lite методом LoRA~\cite{hu2022lora} через API Yandex Cloud~\cite{yandexgpt2024docs}. На выходе --- дообученная модель с уникальным идентификатором чекпоинта.
    \item \textbf{Инференс на бенчмарке.} Последовательный прогон каждого примера из бенчмарка через модель с фиксированными параметрами генерации. Результат --- XML-разметка с вызовом инструмента.
    \item \textbf{Оценка и конвертация.} Парсинг ответов модели, извлечение имени вызванного инструмента из тегов \texttt{<tool\_call>}, сопоставление с эталонной меткой. Вычисление метрик accuracy, precision, recall, F1 для каждого класса.
    \item \textbf{Статистический анализ.} Применение непараметрических критериев (Cochran Q, McNemar) и бутстрап-оценок доверительных интервалов для проверки статистической значимости различий между моделями.
\end{enumerate}

\subsection{Описание бенчмарков} \label{subsec:benchmarks}

Для оценки качества классификации инструментов используются два бенчмарка, различающихся по составу и назначению (Табл.~\ref{tab:benchmarks}).

\begin{table}[H]
\centering
\caption{Характеристики оценочных бенчмарков.}
\label{tab:benchmarks}
\def\arraystretch{1.2}
\setlength\tabcolsep{0.4em}
\begin{tabular}{lcccc}
    \toprule
    Бенчмарк & Объём & GEO & TAG & Тип данных \\
    \midrule
    \texttt{combined} & 500 & 326 (65,2\%) & 174 (34,8\%) & Реал.\,+ синтет. \\
    \texttt{real}     & 250 & 203 (81,2\%) & 47 (18,8\%)  & Только реальные \\
    \bottomrule
\end{tabular}
\end{table}

\textbf{benchmark-combined} (500 примеров) содержит смесь реальных пользовательских запросов и синтетических примеров. Соотношение классов GEO/TAG составляет 326/174, что отражает естественное преобладание географических запросов в предметной области. Данный бенчмарк предназначен для оценки общей способности модели к классификации инструментов.

\textbf{benchmark-real} (250 примеров) состоит исключительно из реальных пользовательских запросов, собранных из логов AI-агента. Распределение классов существенно смещено: 203 запроса GEO (81,2\%) против 47 запросов TAG (18,8\%). Данный бенчмарк предназначен для оценки качества модели в условиях, максимально приближённых к промышленной эксплуатации.

\subsubsection{Политика разметки пограничных случаев}

При создании бенчмарков применялась единая политика разметки, определяющая отнесение запроса к классу GEO или TAG в неоднозначных ситуациях:
\begin{itemize}
    \item \textbf{GEO}: запрос содержит конкретные имена собственные --- названия станций метро (<<ст. м. Сокол>>), улиц (<<ул. Тверская>>), районов, городов, конкретных адресов. Критерий --- наличие именованной географической сущности, требующей разрешения через геокодер.
    \item \textbf{TAG}: запрос описывает класс инфраструктуры или характеристику объекта без привязки к конкретному имени собственному. Примеры: <<рядом с парком>>, <<с камином>>, <<у метро>> (как категория транспортной доступности, без указания конкретной станции).
\end{itemize}

\noindent Ключевое разграничение: выражение <<у метро Сокол>> классифицируется как GEO (содержит имя собственное), тогда как <<у метро>> --- как TAG (описывает характеристику без конкретной привязки).

\subsection{Обучающие датасеты} \label{subsec:datasets}

Для исследования влияния объёма и типа обучающих данных на качество классификации были подготовлены шесть основных обучающих датасетов (Табл.~\ref{tab:datasets}).

\begin{table}[H]
\centering
\caption{Характеристики обучающих датасетов.}
\label{tab:datasets}
\def\arraystretch{1.2}
\setlength\tabcolsep{0.35em}
\begin{tabular}{lclccl}
    \toprule
    Датасет & Объём & Тип & GEO & TAG & Гуман. \\
    \midrule
    \texttt{synth-50}   & 50  & Синтет.       & 54,0\% & 46,0\% & Да \\
    \texttt{synth-250}  & 250 & Синтет.       & 53,6\% & 46,4\% & Да \\
    \texttt{synth-500}  & 500 & Синтет.       & 55,8\% & 44,2\% & Да \\
    \texttt{synth-1000} & 908 & Синтет.       & 52,0\% & 48,0\% & Да \\
    \texttt{real-250}   & 250 & Реал.         & 82,4\% & 17,6\% & Нет \\
    \texttt{mixed-500}  & 500 & Реал.\,+ синт. & 68,0\% & 32,0\% & Нет \\
    \bottomrule
\end{tabular}
\end{table}

Синтетические датасеты (\texttt{train-synthetic-*}) сгенерированы по методологии Self-Instruct~\cite{wang2023selfinstruct} и прошли процедуру гуманизации --- автоматическое преобразование, приближающее стилистику и лексику запросов к реальным пользовательским формулировкам. Классовый баланс в синтетических датасетах поддерживается на уровне, близком к 50/50.

Датасет \texttt{train-real-250} содержит только реальные пользовательские запросы с естественным распределением классов (82,4\% GEO / 17,6\% TAG), что отражает доминирование географических запросов в предметной области.

Датасет \texttt{train-mixed-500} представляет собой объединение реальных и синтетических примеров, что позволяет исследовать эффект смешивания типов данных при обучении.

Выбор такого набора датасетов мотивирован необходимостью ответить на следующие исследовательские вопросы: (а) как объём обучающей выборки влияет на качество (серия \texttt{synthetic-50/250/500/1000}); (б) превосходят ли реальные данные синтетические при равном объёме (\texttt{real-250} vs \texttt{synthetic-250}); (в) даёт ли смешивание типов данных преимущество (\texttt{mixed-500} vs \texttt{synthetic-500}).

Помимо основных шести датасетов, были обучены дополнительные версии моделей \texttt{lite\_latest} (v1, v2, v3) и \texttt{lite\_rc} (v1, v2) на корпусе из приблизительно 1500 синтетических примеров. Эти модели представляют собой результаты итеративного улучшения конфигурации обучения.

\subsection{Протокол дообучения} \label{subsec:training}

\subsubsection{Метод дообучения}

Дообучение выполняется методом LoRA~\cite{hu2022lora} (раздел~\ref{subsec:lora}) через API Yandex Cloud~\cite{yandexgpt2024docs}. Yandex Cloud API предоставляет дообучение как управляемый сервис: пользователь загружает обучающий датасет в формате JSONL, указывает базовую модель и запускает задачу обучения. Внутренние гиперпараметры LoRA (ранг $r$, коэффициент масштабирования $\alpha$, learning rate) управляются платформой и недоступны для прямого задания пользователем.

\subsubsection{Базовая модель}

В качестве базовой модели для дообучения используется YandexGPT Lite --- компактная версия семейства YandexGPT, оптимизированная для задач с низкой латентностью. Для сравнения в роли baseline выступает более крупная модель YandexGPT~5.1~Pro (\texttt{pro\_baseline}), используемая без дообучения.

\subsubsection{Конфигурация инференса}

При проведении инференса на бенчмарках для всех моделей применяются единые параметры генерации:
\begin{itemize}
    \item \textbf{Temperature}: $0{,}1$ --- низкое значение для обеспечения детерминированности и воспроизводимости ответов;
    \item \textbf{Max tokens}: $4096$ --- верхняя граница длины генерируемого ответа;
    \item \textbf{Парсинг}: извлечение имени вызванного инструмента из XML-тегов \texttt{<tool\_call>}.
\end{itemize}

\subsection{Дизайн эксперимента} \label{subsec:design}

\subsubsection{Перечень моделей}

В эксперименте участвуют 12 моделей: 1 baseline-модель без дообучения и 11 дообученных вариантов YandexGPT Lite. Полный перечень приведён в Табл.~\ref{tab:models}.

\begin{table}[H]
\centering
\caption{Перечень моделей, участвующих в эксперименте.}
\label{tab:models}
\def\arraystretch{1.2}
\setlength\tabcolsep{0.35em}
\begin{tabular}{llll}
    \toprule
    \textbf{Модель} & \textbf{База} & \textbf{Обучающий датасет} & \textbf{Объём} \\
    \midrule
    \texttt{pro\_baseline}        & GPT~5.1~Pro & --- (без дообучения) & --- \\
    \texttt{train\_synth\_50}     & GPT~Lite    & \texttt{synth\nobreakdash-50}   & 50 \\
    \texttt{train\_synth\_250}    & GPT~Lite    & \texttt{synth\nobreakdash-250}  & 250 \\
    \texttt{train\_synth\_500}    & GPT~Lite    & \texttt{synth\nobreakdash-500}  & 500 \\
    \texttt{train\_synth\_1000}   & GPT~Lite    & \texttt{synth\nobreakdash-1000} & 908 \\
    \texttt{train\_real\_250}     & GPT~Lite    & \texttt{real\nobreakdash-250}       & 250 \\
    \texttt{train\_mixed\_500}    & GPT~Lite    & \texttt{mixed\nobreakdash-500}      & 500 \\
    \texttt{lite\_latest\_v1}     & GPT~Lite    & ${\sim}$1500 синтет.     & ${\sim}$1500 \\
    \texttt{lite\_latest\_v2}     & GPT~Lite    & ${\sim}$1500 синтет.     & ${\sim}$1500 \\
    \texttt{lite\_latest\_v3}     & GPT~Lite    & ${\sim}$1500 синтет.     & ${\sim}$1500 \\
    \texttt{lite\_rc\_v1}         & GPT~Lite    & ${\sim}$1500 синтет.     & ${\sim}$1500 \\
    \texttt{lite\_rc\_v2}         & GPT~Lite    & ${\sim}$1500 синтет.     & ${\sim}$1500 \\
    \bottomrule
\end{tabular}
\end{table}

\subsubsection{Матрица экспериментов}

Каждая из 12 конфигураций моделей оценивается на каждом из 2 бенчмарков, что даёт $12 \times 2 = 24$ цикла оценки (evaluation runs). Для каждого цикла фиксируются бинарные исходы (верно/неверно) по каждому примеру бенчмарка, что обеспечивает возможность применения попарных статистических тестов на связанных выборках.

\subsubsection{Статистические тесты} \label{subsubsec:stat_tests}

Для анализа результатов применяется трёхуровневая схема статистического тестирования (формальные определения всех критериев приведены в разделе~\ref{subsec:statistical-tests}).

\textbf{Глобальный тест.} Критерий Кохрена~$Q$~\cite{cochran1950comparison} проверяет нулевую гипотезу о равенстве вероятностей правильной классификации для всех $k$ моделей. При отклонении $H_0$ выполняются попарные сравнения.

\textbf{Попарные сравнения.} Для каждой пары моделей применяется точный критерий Макнемара~\cite{mcnemar1947note}. При $\binom{12}{2} = 66$ парах применяется поправка Бонферрони~\cite{bonferroni1936teoria}: $\alpha_{\mathrm{adj}} = 0{,}05 / 66 \approx 0{,}0008$.

\textbf{Доверительные интервалы.} Для оценки accuracy каждой модели строятся 95\%-доверительные интервалы методом бутстрапа~\cite{efron1979bootstrap} ($B = 1000$, $\mathrm{seed} = 42$).

\subsubsection{Контроль воспроизводимости}

Для обеспечения воспроизводимости результатов приняты следующие меры:
\begin{itemize}
    \item Фиксированное значение temperature $= 0{,}1$ при инференсе для минимизации стохастичности генерации.
    \item Фиксированное зерно $\mathrm{seed} = 42$ для всех процедур бутстрапа.
    \item Сохранение уникальных идентификаторов чекпоинтов всех дообученных моделей в Yandex Cloud.
    \item Сохранение полных логов инференса (запрос, ответ модели, извлечённая метка) для каждого эксперимента.
    \item Использование единых бенчмарков без перемешивания или ресемплирования между запусками.
\end{itemize}
